% Chinese setup.
% Overleaf has ctex and the Fandol fonts, but not this machine's Droid font
% or the absolute /data/workspace paths. Use ctex there, and the bundled
% TrueType file only when ctex is absent.
\IfFileExists{ctex.sty}{%
  \usepackage[fontset=fandol]{ctex}%
  \newcommand{\cjkfont}{}%
  \newcommand{\cjkbr}{}%
  \newcount\cjkbreakoff
}{%
  \directlua{tex.enableprimitives("", tex.extraprimitives("pdftex", "luatex"))}%
  \usepackage{fontspec}%
  \defaultfontfeatures{Ligatures=TeX}%
  \IfFileExists{fonts/DroidSansFallback.ttf}{%
    \newfontfamily\cjkfont{DroidSansFallback}[
      Path=fonts/,
      Extension=.ttf,
      UprightFont=*,
      BoldFont=*,
      BoldFeatures={FakeBold=2.4},
      ItalicFont=*,
      ItalicFeatures={FakeSlant=0.16},
      BoldItalicFont=*,
      BoldItalicFeatures={FakeBold=2.4,FakeSlant=0.16}
    ]%
  }{%
    \newfontfamily\cjkfont{DroidSansFallback}[
      Path=../../fonts/,
      Extension=.ttf,
      UprightFont=*,
      BoldFont=*,
      BoldFeatures={FakeBold=2.4},
      ItalicFont=*,
      ItalicFeatures={FakeSlant=0.16},
      BoldItalicFont=*,
      BoldItalicFeatures={FakeBold=2.4,FakeSlant=0.16}
    ]%
  }%
  \DeclareRobustCommand{\cjkbr}{\ifmmode\else\hskip 0pt plus 0.15pt minus 0.05pt\relax\fi}%
  \newcount\cjkbreakoff
  \directlua{%
    local name = "cjkbreak.lua"
    local path = kpse.find_file(name, "tex") or kpse.find_file(name) or name
    dofile(path)
  }%
}
